\documentclass[12pt,a4paper]{article}

% ========== Packages ==========
\usepackage[margin=2.5cm]{geometry}
\usepackage{graphicx}
\usepackage{booktabs}
\usepackage{multirow}
\usepackage{amsmath}
\usepackage{amssymb}
\usepackage{hyperref}
\usepackage{enumitem}
\usepackage{caption}
\usepackage{subcaption}
\usepackage{float}
\usepackage{listings}
\usepackage{xcolor}
\usepackage{fancyhdr}
\usepackage{titlesec}

% ========== Page Style ==========
\pagestyle{fancy}
\fancyhf{}
\fancyhead[C]{\small Wind Turbine Blade Defect Detection Based on Improved YOLOv11}
\fancyfoot[C]{\thepage}
\renewcommand{\headrulewidth}{0.4pt}

% ========== Code Style ==========
\lstset{
  basicstyle=\small\ttfamily,
  backgroundcolor=\color{gray!5},
  frame=single,
  breaklines=true,
  columns=flexible,
  numbers=left,
  numberstyle=\tiny\color{gray},
  keywordstyle=\color{blue},
  commentstyle=\color{green!60!black},
  stringstyle=\color{red!80!black},
}

% ========== Figure Path ==========
\graphicspath{{../data/docs/figures/}}

% ========== Title ==========
\title{\textbf{Wind Turbine Blade Defect Detection System\\Based on Improved YOLOv11}}
\author{Course Design Report}
\date{\today}

\begin{document}

\maketitle

% ============================================================
% Abstract
% ============================================================
\begin{abstract}
Wind energy is a cornerstone of clean energy transition, and defect detection of wind turbine blades is critical for ensuring safe operation. Traditional manual inspection methods are inefficient and subjective, making them inadequate for large-scale wind farm maintenance. This paper presents an automatic defect detection system for wind turbine blade surface defects (cracks and erosion) based on YOLO-series object detection models.

First, a wind turbine blade defect detection dataset containing 1,096 images with two defect categories (crack and erosion) was constructed in YOLO format. Second, YOLOv11n was established as the baseline model, fine-tuned from COCO pretrained weights, achieving a baseline mAP@0.5 of 78.01\%. Third, systematic comparison experiments were designed to evaluate YOLOv5n, YOLOv8n, and YOLOv11n under identical conditions. Finally, ablation studies on freeze strategies were conducted to investigate the impact of different frozen layer counts on transfer learning effectiveness for small datasets.

Experimental results demonstrate that YOLOv8n achieves the best performance on this dataset with mAP@0.5 of 82.80\%, a 4.79 percentage point improvement over the baseline. Full fine-tuning (freeze=0) of YOLOv11n (80.35\%) outperforms all frozen strategies, confirming that COCO pretrained weights are essential for small dataset detection tasks.

\medskip
\noindent\textbf{Keywords}: Wind turbine blade; Defect detection; YOLOv11; Transfer learning; Freeze strategy
\end{abstract}

\newpage
\tableofcontents
\newpage

% ============================================================
% Chapter 1: Introduction
% ============================================================
\section{Introduction}

\subsection{Research Background and Significance}

Wind power generation is a vital pillar of the global energy transition. As of 2025, global installed wind power capacity has exceeded 1,000\,GW, with China being the world's largest wind power market. Wind turbine blades, as core components, are continuously exposed to complex natural environments and are susceptible to surface defects such as cracks, corrosion, and lightning damage. Undetected defects may lead to catastrophic blade failures, causing significant economic losses and safety risks.

Traditional blade inspection relies primarily on manual visual examination, which suffers from several limitations:
\begin{itemize}[noitemsep]
  \item \textbf{Low efficiency}: Inspecting a single turbine's blades takes several hours, with full wind farm inspections spanning months;
  \item \textbf{Subjectivity}: Detection results depend on inspector experience, leading to high miss and false alarm rates;
  \textbf{Safety risks}: High-altitude work poses hazards, and inspections cannot be performed in adverse weather;
  \item \textbf{High costs}: Professional inspectors and specialized equipment are required, driving up maintenance costs.
\end{itemize}

Therefore, developing automated and intelligent wind turbine blade defect detection systems has significant engineering value. The advancement of deep learning-based object detection provides new solutions to this problem.

\subsection{Related Work}

\subsubsection{Traditional Image Processing Methods}
Early defect detection research relied on traditional image processing techniques, including edge detection (Canny, Sobel), morphological operations, and threshold segmentation. These methods are effective for simple defects but lack robustness against complex backgrounds and varying illumination conditions.

\subsubsection{Deep Learning Object Detection Methods}
In recent years, deep learning-based object detection methods have made significant progress in industrial defect detection:

\textbf{Two-stage methods}: Represented by Faster R-CNN, these generate region proposals first then classify and localize, achieving high accuracy but slower speed.

\textbf{Single-stage methods}: Represented by the YOLO (You Only Look Once) series, these treat detection as a regression problem, predicting bounding boxes and class probabilities directly from image pixels in a single forward pass, achieving a good balance between speed and accuracy.

\subsubsection{Attention Mechanisms and Feature Fusion}
Attention mechanisms (SE, CBAM, CA) enhance important features through adaptive weighting and are widely used in defect detection. Feature fusion techniques (FPN, PANet, BiFPN) improve small object detection through multi-scale feature aggregation.

\subsection{Contributions}

The main contributions of this paper are:
\begin{enumerate}[noitemsep]
  \item Constructed a wind turbine blade defect detection dataset containing 1,096 images with two defect categories (crack and erosion);
  \item Established YOLOv11n as the baseline model, fine-tuned from COCO pretrained weights, achieving mAP@0.5 of 78.01\%;
  \item Designed systematic comparison experiments evaluating YOLOv5n, YOLOv8n, and YOLOv11n under identical conditions;
  \item Conducted freeze strategy ablation studies (freeze=0/5/10/11) to investigate transfer learning effectiveness on small datasets;
  \item Key finding: YOLOv8n achieves the best performance on small datasets (mAP@0.5 = 82.80\%), and COCO pretrained weights are critical for small dataset detection.
\end{enumerate}

% ============================================================
% Chapter 2: Related Technologies
% ============================================================
\section{Related Technologies}

\subsection{YOLO Object Detection Framework}

YOLO (You Only Look Once) is one of the most popular single-stage object detection algorithms. Its core idea is to transform object detection into an end-to-end regression problem, predicting bounding box locations and class probabilities simultaneously through a single forward pass.

A YOLO model typically consists of three parts:
\begin{itemize}[noitemsep]
  \item \textbf{Backbone}: Extracts multi-level features from input images (e.g., CSPDarknet53, C2f modules);
  \item \textbf{Neck}: Performs multi-scale feature fusion (e.g., PANet, BiFPN);
  \item \textbf{Head}: Performs final bounding box regression and classification.
\end{itemize}

The YOLO versions involved in this study include:
\begin{itemize}[noitemsep]
  \item \textbf{YOLOv5n}: CSPDarknet53 Backbone + PANet Neck, 2.50M parameters, 7.1 GFLOPs;
  \item \textbf{YOLOv8n}: Improved Backbone + C2f module + PANet Neck, 3.01M parameters, 8.1 GFLOPs;
  \item \textbf{YOLOv11n}: C3k2 + C2PSA attention + SPPF, 2.62M parameters, 6.3 GFLOPs.
\end{itemize}

\subsection{Attention Mechanisms}

Attention mechanisms enhance feature representation by adaptively weighting important features.

\subsubsection{Channel Attention}
The SE (Squeeze-and-Excitation) module obtains channel statistics via global average pooling, then generates channel attention weights through two fully connected layers.

\subsubsection{Spatial Attention}
Spatial attention focuses on the importance of different spatial positions in images, typically generating spatial attention maps through convolutional operations.

\subsubsection{Coordinate Attention (CA)}
Coordinate Attention decomposes channel attention into two one-dimensional feature encodings, capturing horizontal and vertical positional information while preserving precise location awareness.

\subsubsection{PSA Mechanism}
PSA (Pixel Shift Attention) is introduced in YOLOv11, modeling pixel relationships within local windows through pixel shift operations. This paper uses the C2PSA module with the repeat parameter increased from 2 to 4 for enhanced attention.

\subsection{Feature Fusion Techniques}

Multi-scale feature fusion is key to improving detection performance:
\begin{itemize}[noitemsep]
  \item \textbf{FPN} (Feature Pyramid Network): Top-down feature pyramid传递高层语义特征;
  \item \textbf{PANet} (Path Aggregation Network): Adds bottom-up path aggregation on top of FPN;
  \item \textbf{BiFPN} (Bidirectional Feature Pyramid Network): Bidirectional feature pyramid for efficient multi-scale fusion.
\end{itemize}

\subsection{Lightweight Techniques}

For edge deployment, lightweight techniques reduce model parameters and computation:
\begin{itemize}[noitemsep]
  \item \textbf{Depthwise separable convolution}: Decomposes standard convolution into depthwise and pointwise operations;
  \item \textbf{GhostNet/GSConv}: Uses lightweight feature generation modules;
  \item \textbf{Reducing network depth}: E.g., reducing C3k2 repeat from 2 to 1.
\end{itemize}

% ============================================================
% Chapter 3: Method Design
% ============================================================
\section{Method Design}

\subsection{Dataset and Preprocessing}

\subsubsection{Data Source}
The dataset consists of actual wind turbine blade inspection images.

\subsubsection{Defect Categories}
The dataset contains two defect types:
\begin{itemize}[noitemsep]
  \item \textbf{Crack}: Linear fractures on blade surfaces, caused by fatigue loading or manufacturing defects;
  \item \textbf{Erosion}: Material degradation on blade surfaces, typically caused by environmental factors (rain, salt spray).
\end{itemize}

\subsubsection{Dataset Scale}
The dataset contains 1,096 images, split in a 7:1.5:1.5 ratio:
\begin{itemize}[noitemsep]
  \item Training set: 764 images
  \item Validation set: 166 images
  \item Test set: 166 images
\end{itemize}

\subsubsection{Data Format}
YOLO format annotations are used, with each image paired with a txt file containing class IDs and normalized bounding box coordinates ($x_{center}$, $y_{center}$, $width$, $height$).

\subsubsection{Data Augmentation}
Multiple augmentation strategies are applied during training:
\begin{itemize}[noitemsep]
  \item Mosaic augmentation ($m=1.0$): Combines 4 images into one, enriching target scales and backgrounds;
  \item MixUp augmentation ($m=0.1$): Randomly mixes two images to enhance robustness;
  \item CopyPaste augmentation ($m=0.1$): Copies and pastes target regions to increase positive samples.
\end{itemize}

\subsection{Baseline Model Selection}

YOLOv11n was selected as the baseline model for the following reasons:
\begin{itemize}[noitemsep]
  \item Lightweight design with only 2.62M parameters and 6.3 GFLOPs, suitable for resource-constrained deployment;
  \item Introduces C2PSA attention mechanism, enhancing feature representation while maintaining efficiency;
  \item As the latest YOLO version, it represents the current state-of-the-art in single-stage detectors.
\end{itemize}

The baseline was fine-tuned from COCO pretrained weights (yolo11n.pt), achieving mAP@0.5 = 78.01\%.

\subsection{Comparison Experiment Design}

Three mainstream YOLO lightweight models were compared under identical conditions:

\begin{table}[H]
\centering
\caption{Model Configurations for Comparison Experiments}
\begin{tabular}{lccc}
\toprule
Model & Backbone & Neck & Parameters \\
\midrule
YOLOv5n & CSPDarknet53 & PANet & 2.50M \\
YOLOv8n & C2f-Backbone & PANet & 3.01M \\
YOLOv11n & C3k2+C2PSA & PANet & 2.62M \\
\bottomrule
\end{tabular}
\end{table}

All models were fine-tuned from COCO pretrained weights with unified training configuration:
\begin{itemize}[noitemsep]
  \item Epochs: 150, Batch size: 8, Image size: 640$\times$640
  \item Optimizer: AdamW (auto), Initial learning rate: 0.001
  \item LR schedule: Cosine annealing (lr0=0.001 $\to$ lrf=0.0001)
  \item Warmup: 5 epochs, warmup momentum: 0.8
  \item Early stopping: patience=30
\end{itemize}

\subsection{Ablation Study Design --- Freeze Strategy}

Different frozen layer counts were compared on YOLOv11n to investigate transfer learning effectiveness:
\begin{itemize}[noitemsep]
  \item \textbf{freeze=0}: No freezing, full fine-tuning (from COCO pretrained)
  \item \textbf{freeze=5}: Freeze first 5 layers (early Backbone)
  \item \textbf{freeze=10}: Freeze first 10 layers (nearly all Backbone)
  \item \textbf{freeze=11}: Freeze all Backbone, train only Head
\end{itemize}

\subsection{Transfer Learning Strategy}

Key finding: Models trained from random initialization (built from YAML) perform extremely poorly (mAP50 of only 0.43--0.57), demonstrating that small datasets (764 training images) are insufficient for training detection models from scratch.

Solution: Transfer learning from COCO pretrained weights or baseline best.pt for fine-tuning.

% ============================================================
% Chapter 4: Experiments and Analysis
% ============================================================
\section{Experiments and Analysis}

\subsection{Experimental Environment}

\begin{table}[H]
\centering
\caption{Experimental Environment}
\begin{tabular}{ll}
\toprule
Item & Configuration \\
\midrule
GPU & NVIDIA RTX 4060 Laptop (8GB VRAM) \\
CUDA & 13.0 \\
PyTorch & 2.6.0 \\
ultralytics & 8.4.48 \\
Python & 3.9 \\
OS & Windows 11 \\
\bottomrule
\end{tabular}
\end{table}

\subsection{Evaluation Metrics}

The following metrics were used:
\begin{itemize}[noitemsep]
  \item \textbf{mAP@0.5}: Mean Average Precision at IoU=0.5, primary metric;
  \item \textbf{mAP@0.5:0.95}: Mean AP across multiple IoU thresholds, stricter evaluation;
  \item \textbf{Precision}: Proportion of correct detections among all detections;
  \item \textbf{Recall}: Proportion of correctly detected targets among all ground truths.
\end{itemize}

Model parameters (Params) and computational cost (GFLOPs) were also recorded.

\subsection{Model Comparison Experiments}

\begin{table}[H]
\centering
\caption{Comparison Experiment Results}
\label{tab:comparison}
\begin{tabular}{lccccc}
\toprule
Model & mAP@0.5 & mAP@0.5:0.95 & Precision & Recall & Params \\
\midrule
YOLOv11n (Baseline) & 0.7801 & 0.4767 & 0.8337 & 0.6769 & 2.62M \\
YOLOv5n & 0.8078 & 0.4887 & 0.8581 & 0.7000 & 2.50M \\
\textbf{YOLOv8n} & \textbf{0.8280} & \textbf{0.5218} & \textbf{0.9024} & \textbf{0.7237} & 3.01M \\
YOLOv11n+Freeze & 0.7644 & 0.4629 & 0.8488 & 0.6667 & 2.62M \\
\bottomrule
\end{tabular}
\end{table}

\textbf{Analysis}:

\begin{enumerate}[noitemsep]
  \item \textbf{YOLOv8n achieves the best performance}: mAP@0.5 of 82.80\%, a 4.79 percentage point improvement over the baseline, with precision reaching 90.24\%. The C2f module in YOLOv8n demonstrates superior feature extraction and representation capabilities for small dataset defect detection.

  \item \textbf{YOLOv5n outperforms the baseline}: Despite similar parameter counts (2.50M vs 2.62M), YOLOv5n achieves higher mAP@0.5 (80.78\% vs 78.01\%), indicating that the classic CSPDarknet53 backbone has stronger feature extraction capability on this dataset compared to YOLOv11's improved backbone.

  \item \textbf{YOLOv11n attention mechanism shows limited advantage}: The C2PSA attention mechanism does not fully leverage its feature selection capability on the small dataset (764 training images), possibly because insufficient data prevents effective attention weight learning.

  \item \textbf{Freezing backbone degrades performance}: YOLOv11n+Freeze (freeze=11) achieves mAP@0.5 of 76.44\%, lower than the full-training baseline (78.01\%), though precision improves (84.88\% vs 83.37\%), suggesting freeze strategies reduce overfitting but also limit learning capacity.
\end{enumerate}

\subsection{Ablation Study --- Freeze Strategy}

\begin{table}[H]
\centering
\caption{Freeze Strategy Ablation Results}
\label{tab:ablation}
\begin{tabular}{lccccc}
\toprule
Strategy & Frozen Layers & mAP@0.5 & vs Baseline & Precision & Recall \\
\midrule
No freeze & 0 & 0.8035 & +2.97\% & 0.8492 & 0.7232 \\
Partial freeze & 5 & 0.6582 & $-$15.63\% & 0.7689 & 0.5315 \\
Most freeze & 10 & 0.7479 & $-$4.13\% & 0.7669 & 0.6575 \\
Full freeze & 11 & 0.7644 & $-$2.02\% & 0.8488 & 0.6667 \\
\bottomrule
\end{tabular}
\end{table}

\textbf{Analysis}:

\begin{enumerate}[noitemsep]
  \item \textbf{Full fine-tuning achieves the best result}: freeze=0 (no freezing) achieves mAP@0.5 of 80.35\%, a 2.97\% improvement over the baseline. Fine-tuning all layers from COCO pretrained weights with a low learning rate leverages pretrained features while fully adapting to the target dataset.

  \item \textbf{freeze=5 performs the worst}: mAP@0.5 of only 65.82\%, far below other strategies. Freezing the first 5 layers (early low-level feature layers) may prevent the network from learning dataset-specific low-level feature representations, and fine-tuning only the later layers is insufficient to compensate.

  \item \textbf{freeze=10 and freeze=11 show moderate performance}: Both achieve similar results (74.79\% vs 76.44\%). Freezing most of the Backbone and fine-tuning only Neck and Head limits model capacity but reduces overfitting.

  \item \textbf{Non-monotonic relationship}: The relationship between frozen layer count and detection performance is not simply monotonic. freeze=5 performs worst, suggesting that partial freezing may cause incoordination between feature extraction and task adaptation.
\end{enumerate}

% ============================================================
% Chapter 5: Conclusion
% ============================================================
\section{Conclusion and Future Work}

\subsection{Summary}

This paper addresses the wind turbine blade defect detection problem through the following work:
\begin{enumerate}[noitemsep]
  \item Constructed a wind turbine blade defect detection dataset containing 1,096 images with two defect categories;
  \item Established YOLOv11n as the baseline model with mAP@0.5 of 78.01\%;
  \item Designed systematic comparison experiments evaluating YOLOv5n, YOLOv8n, and YOLOv11n;
  \item Conducted freeze strategy ablation studies (freeze=0/5/10/11) to investigate transfer learning effectiveness.
\end{enumerate}

\subsection{Key Findings}

\begin{enumerate}[noitemsep]
  \item \textbf{YOLOv8n is the optimal model}: Achieves mAP@0.5 of 82.80\% on the wind turbine blade defect detection task, the highest among all compared models.

  \item \textbf{Full fine-tuning outperforms freeze strategies}: Full fine-tuning from COCO pretrained weights (freeze=0) achieves 80.35\%, surpassing all freeze strategies. Low learning rate full fine-tuning is an effective transfer learning strategy for small datasets.

  \item \textbf{Freeze strategies have limited effectiveness}: Freezing the backbone (freeze=11) actually degrades performance (76.44\%), and freeze=5 performs worst (65.82\%), indicating that freeze strategies limit model learning capacity on small datasets.

  \item \textbf{COCO pretrained weights are critical}: Random initialization training from YAML achieves mAP50 of only 0.43--0.57, far below transfer learning results, confirming that pretrained weights are indispensable for small dataset detection.

  \item \textbf{YOLOv11 attention mechanism shows limited advantage on small datasets}: The C2PSA attention mechanism does not fully leverage its capability on 764 training images, possibly requiring larger datasets to demonstrate its benefits.
\end{enumerate}

\subsection{Limitations and Future Work}

\begin{enumerate}[noitemsep]
  \item \textbf{Limited data}: Only 764 training images; expanding with more defect categories and data augmentation strategies (rotation, flipping) could improve performance;
  \item \textbf{No complex architectural improvements explored}: BiFPN feature fusion, GhostNet lightweight modules, and CA attention mechanisms remain unexplored;
  \item \textbf{No real-time optimization considered}: TensorRT export, model quantization (INT8/FP16), and pruning for deployment optimization;
  \item \textbf{Stronger pretrained models}: YOLOv11s/m or self-supervised pretraining strategies could enhance performance;
  \item \textbf{Extension to semantic segmentation}: YOLOv11-seg could provide pixel-level defect localization for more precise boundary information.
\end{enumerate}

% ============================================================
% References
% ============================================================
\section*{References}
\addcontentsline{toc}{section}{References}

\begin{enumerate}[label={[\arabic*]}, noitemsep]
  \item Redmon J, Divvala S, Girshick R, et al. You only look once: Unified, real-time object detection[C]//Proceedings of the IEEE conference on computer vision and pattern recognition. 2016: 779--788.
  \item Jocher G, Chaurasia A, Qiu J. Ultralytics YOLOv8[J]. GitHub repository, 2023.
  \item Li C, Li L, Geng Y, et al. YOLOv11: An improved real-time object detection algorithm[J]. arXiv preprint arXiv:2410.17725, 2024.
  \item Ren S, He K, Girshick R, et al. Faster R-CNN: Towards real-time object detection with region proposal networks[J]. IEEE transactions on pattern analysis and machine intelligence, 2017, 39(6): 1137--1149.
  \item Lin T Y, Dollár P, Girshick R, et al. Feature pyramid networks for object detection[C]//Proceedings of the IEEE conference on computer vision and pattern recognition. 2017: 2117--2125.
  \item Liu S, Qi L, Qin H, et al. Path aggregation network for instance segmentation[C]//Proceedings of the IEEE conference on computer vision and pattern recognition. 2018: 8759--8768.
  \item Tan M, Pang R, Le Q V. Efficientdet: Scalable and efficient object detection[C]//Proceedings of the IEEE/CVF conference on computer vision and pattern recognition. 2020: 10781--10790.
  \item Hu J, Shen L, Sun G. Squeeze-and-excitation networks[C]//Proceedings of the IEEE conference on computer vision and pattern recognition. 2018: 7132--7141.
  \item Woo S, Park J, Lee J Y, et al. CBAM: Convolutional block attention module[C]//Proceedings of the European conference on computer vision. 2018: 3--19.
  \item Hou Q, Zhou D, Feng J. Coordinate attention for efficient mobile network design[C]//Proceedings of the IEEE/CVF conference on computer vision and pattern recognition. 2021: 13713--13722.
  \item Jocher G, Chaurasia A, Qiu J. Ultralytics YOLOv5[J]. GitHub repository, 2020.
  \item Wang C Y, Bochkovskiy A, Liao H Y M. YOLOv7: Trainable bag-of-freebies sets new state-of-the-art for real-time object detectors[C]//Proceedings of the IEEE/CVF conference on computer vision and pattern recognition. 2023: 7464--7475.
\end{enumerate}

% ============================================================
% Appendix
% ============================================================
\newpage
\appendix
\section{Full Training Configuration}

\begin{verbatim}
# YOLOv11n Baseline Training Configuration
model: yolo11n.pt  # COCO pretrained
data: wind_turbine_2cls.yaml
epochs: 150
batch: 8
imgsz: 640
optimizer: auto  # AdamW
lr0: 0.001
lrf: 0.0001
cos_lr: True
warmup_epochs: 5
warmup_momentum: 0.8
weight_decay: 0.0005
patience: 30
mosaic: 1.0
mixup: 0.1
copy_paste: 0.1
device: 0
workers: 4
\end{verbatim}

\section{Model Architecture Details}

YOLOv11n network structure contains 101 layers:
\begin{itemize}[noitemsep]
  \item Backbone: Conv + C3k2 + SPPF + C2PSA
  \item Neck: Upsample + Concat + C3k2 (PANet structure)
  \item Head: Detect (3 detection scales: P3/P4/P5)
  \item Parameters: 2.62M, GFLOPs: 6.3
\end{itemize}

\end{document}
